% TOP_PROBLEM: {"id": 1444, "title": "Tutte's 4-Flow Conjecture for Petersen-Minor-Free Graphs", "category_id": 3, "category": {"id": 3, "name": "graph_theory", "display_name": "Graph Theory", "description": "Problems involving graphs, networks, and their properties.", "slug": "graph-theory", "order_index": 3, "created_at": "2026-07-31T15:26:25.671Z"}, "status": "open"}
% FIRST_POSTED: null
% ENTRY_KIND: statement_only
% Imported from: batch07.tex; UnsolvedMath ID 1444
\documentclass[11pt]{article}
\usepackage[margin=25mm]{geometry}
\usepackage{fontspec}
\setmainfont{Latin Modern Roman}
\usepackage{amsmath,amssymb,mathtools}
\usepackage{unicode-math}
\setmathfont{Latin Modern Math}
\usepackage{xurl,hyperref}
\hypersetup{colorlinks=true,urlcolor=blue}
\setcounter{secnumdepth}{0}
\setlength{\parindent}{0pt}
\setlength{\parskip}{5pt}
\setlength{\emergencystretch}{3em}
\newcommand{\sref}[2]{\hyperlink{source:#1:#2}{[S#2]}}
\newcommand{\cataloguescope}[1]{\paragraph{Recorded target.}#1}
\begin{document}
% UnsolvedMath ID 1444; source catalogue code GRAPH-057
\setcounter{section}{1443}
\section{Tutte's 4-Flow Conjecture for Petersen-Minor-Free Graphs}\label{1444}
{\small\sffamily UnsolvedMath ID 1444 · Graph Theory}
\subsection{Definitions and mathematical statement}

\paragraph{Shared notation.}$\mathbb N=\{1,2,\ldots\}$, and $\mathbb Z,\mathbb Q,\mathbb R,\mathbb C$ denote the integers, rationals, reals and complex numbers. Any other conventions are specified locally.


Graphs are finite and undirected, with loops and parallel edges permitted. An edge is a bridge if deleting it increases the number of connected components. A graph minor is obtained by deleting vertices or edges and contracting non-loop edges. Define the Petersen graph $P$ using vertices $a_i,b_i$ for $i\in\mathbb Z/5\mathbb Z$ and edges $a_ia_{i+1}$, $a_ib_i$, and $b_ib_{i+2}$. A graph is Petersen-minor-free if it has no minor isomorphic to $P$. A nowhere-zero $4$-flow consists of an orientation of the edges and an integer assignment $f:E(G)\to\{-3,-2,-1,1,2,3\}$ satisfying, at every vertex $v$,
\[
 \sum_{e\text{ leaving }v}f(e)=\sum_{e\text{ entering }v}f(e).
\]
A loop contributes once to each sum; the empty assignment is a flow on an edgeless graph.
\paragraph{Statement recorded in the supplied source.}
Does every finite bridgeless Petersen-minor-free graph admit a nowhere-zero $4$-flow? No restriction to cubic graphs, connected graphs, or exclusion of a proper minor of the Petersen graph is imposed. \sref{1444}{2}
\subsection{Short English statement}
Can every graph with no bridge and no Petersen graph minor carry a nonzero integer flow of absolute value less than four on every edge?
\subsection{Sources}
\begin{description}
\item[\hypertarget{source:1444:1}{S1}] ULAM.AI and the UnsolvedMath Contributors, UnsolvedMath: A Curated Collection of Open Mathematics Problems, record 1444 (GRAPH-057). CC BY 4.0. \url{https://huggingface.co/datasets/ulamai/UnsolvedMath}.
\item[\hypertarget{source:1444:2}{S2}] József Pintér, Nowhere-Zero 4-Flows in Graphs Excluding a Proper Minor of the Petersen Graph, arXiv:2607.22267v2 (2026), abstract and Sections 1--2. \url{https://arxiv.org/pdf/2607.22267v2}.
\end{description}
\label{end:1444}
\end{document}
